\IfFileExists{stacks-project.cls}{%
\documentclass{stacks-project}
}{%
\documentclass{amsart}
}

% For dealing with references we use the comment environment
\usepackage{verbatim}
\newenvironment{reference}{\comment}{\endcomment}
%\newenvironment{reference}{}{}
\newenvironment{slogan}{\comment}{\endcomment}
\newenvironment{history}{\comment}{\endcomment}

% For commutative diagrams we use Xy-pic
\usepackage[all]{xy}

% We use 2cell for 2-commutative diagrams.
\xyoption{2cell}
\UseAllTwocells

% We use multicol for the list of chapters between chapters
\usepackage{multicol}

% This is generally recommended for better output
\usepackage{lmodern}
\usepackage[T1]{fontenc}

% For cross-file-references
\usepackage{xr-hyper}

% Package for hypertext links:
\usepackage{hyperref}

% For any local file, say "hello.tex" you want to link to please
% use \externaldocument[hello-]{hello}
\externaldocument[introduction-]{introduction}
\externaldocument[conventions-]{conventions}
\externaldocument[sets-]{sets}
\externaldocument[categories-]{categories}
\externaldocument[topology-]{topology}
\externaldocument[sheaves-]{sheaves}
\externaldocument[sites-]{sites}
\externaldocument[stacks-]{stacks}
\externaldocument[fields-]{fields}
\externaldocument[algebra-]{algebra}
\externaldocument[brauer-]{brauer}
\externaldocument[homology-]{homology}
\externaldocument[derived-]{derived}
\externaldocument[simplicial-]{simplicial}
\externaldocument[more-algebra-]{more-algebra}
\externaldocument[smoothing-]{smoothing}
\externaldocument[modules-]{modules}
\externaldocument[sites-modules-]{sites-modules}
\externaldocument[injectives-]{injectives}
\externaldocument[cohomology-]{cohomology}
\externaldocument[sites-cohomology-]{sites-cohomology}
\externaldocument[dga-]{dga}
\externaldocument[dpa-]{dpa}
\externaldocument[sdga-]{sdga}
\externaldocument[hypercovering-]{hypercovering}
\externaldocument[schemes-]{schemes}
\externaldocument[constructions-]{constructions}
\externaldocument[properties-]{properties}
\externaldocument[morphisms-]{morphisms}
\externaldocument[coherent-]{coherent}
\externaldocument[divisors-]{divisors}
\externaldocument[limits-]{limits}
\externaldocument[varieties-]{varieties}
\externaldocument[topologies-]{topologies}
\externaldocument[descent-]{descent}
\externaldocument[perfect-]{perfect}
\externaldocument[more-morphisms-]{more-morphisms}
\externaldocument[flat-]{flat}
\externaldocument[groupoids-]{groupoids}
\externaldocument[more-groupoids-]{more-groupoids}
\externaldocument[etale-]{etale}
\externaldocument[chow-]{chow}
\externaldocument[intersection-]{intersection}
\externaldocument[pic-]{pic}
\externaldocument[weil-]{weil}
\externaldocument[adequate-]{adequate}
\externaldocument[dualizing-]{dualizing}
\externaldocument[duality-]{duality}
\externaldocument[discriminant-]{discriminant}
\externaldocument[derham-]{derham}
\externaldocument[local-cohomology-]{local-cohomology}
\externaldocument[algebraization-]{algebraization}
\externaldocument[curves-]{curves}
\externaldocument[resolve-]{resolve}
\externaldocument[models-]{models}
\externaldocument[functors-]{functors}
\externaldocument[equiv-]{equiv}
\externaldocument[pione-]{pione}
\externaldocument[etale-cohomology-]{etale-cohomology}
\externaldocument[proetale-]{proetale}
\externaldocument[relative-cycles-]{relative-cycles}
\externaldocument[more-etale-]{more-etale}
\externaldocument[trace-]{trace}
\externaldocument[crystalline-]{crystalline}
\externaldocument[spaces-]{spaces}
\externaldocument[spaces-properties-]{spaces-properties}
\externaldocument[spaces-morphisms-]{spaces-morphisms}
\externaldocument[decent-spaces-]{decent-spaces}
\externaldocument[spaces-cohomology-]{spaces-cohomology}
\externaldocument[spaces-limits-]{spaces-limits}
\externaldocument[spaces-divisors-]{spaces-divisors}
\externaldocument[spaces-over-fields-]{spaces-over-fields}
\externaldocument[spaces-topologies-]{spaces-topologies}
\externaldocument[spaces-descent-]{spaces-descent}
\externaldocument[spaces-perfect-]{spaces-perfect}
\externaldocument[spaces-more-morphisms-]{spaces-more-morphisms}
\externaldocument[spaces-flat-]{spaces-flat}
\externaldocument[spaces-groupoids-]{spaces-groupoids}
\externaldocument[spaces-more-groupoids-]{spaces-more-groupoids}
\externaldocument[bootstrap-]{bootstrap}
\externaldocument[spaces-pushouts-]{spaces-pushouts}
\externaldocument[spaces-chow-]{spaces-chow}
\externaldocument[groupoids-quotients-]{groupoids-quotients}
\externaldocument[spaces-more-cohomology-]{spaces-more-cohomology}
\externaldocument[spaces-simplicial-]{spaces-simplicial}
\externaldocument[spaces-duality-]{spaces-duality}
\externaldocument[formal-spaces-]{formal-spaces}
\externaldocument[restricted-]{restricted}
\externaldocument[spaces-resolve-]{spaces-resolve}
\externaldocument[formal-defos-]{formal-defos}
\externaldocument[defos-]{defos}
\externaldocument[cotangent-]{cotangent}
\externaldocument[examples-defos-]{examples-defos}
\externaldocument[algebraic-]{algebraic}
\externaldocument[examples-stacks-]{examples-stacks}
\externaldocument[stacks-sheaves-]{stacks-sheaves}
\externaldocument[criteria-]{criteria}
\externaldocument[artin-]{artin}
\externaldocument[quot-]{quot}
\externaldocument[stacks-properties-]{stacks-properties}
\externaldocument[stacks-morphisms-]{stacks-morphisms}
\externaldocument[stacks-limits-]{stacks-limits}
\externaldocument[stacks-cohomology-]{stacks-cohomology}
\externaldocument[stacks-perfect-]{stacks-perfect}
\externaldocument[stacks-introduction-]{stacks-introduction}
\externaldocument[stacks-more-morphisms-]{stacks-more-morphisms}
\externaldocument[stacks-geometry-]{stacks-geometry}
\externaldocument[moduli-]{moduli}
\externaldocument[moduli-curves-]{moduli-curves}
\externaldocument[examples-]{examples}
\externaldocument[exercises-]{exercises}
\externaldocument[guide-]{guide}
\externaldocument[desirables-]{desirables}
\externaldocument[coding-]{coding}
\externaldocument[obsolete-]{obsolete}
\externaldocument[fdl-]{fdl}
\externaldocument[index-]{index}

% Theorem environments.
%
\theoremstyle{plain}
\newtheorem{theorem}[subsection]{Theorem}
\newtheorem{proposition}[subsection]{Proposition}
\newtheorem{lemma}[subsection]{Lemma}

\theoremstyle{definition}
\newtheorem{definition}[subsection]{Definition}
\newtheorem{example}[subsection]{Example}
\newtheorem{exercise}[subsection]{Exercise}
\newtheorem{situation}[subsection]{Situation}

\theoremstyle{remark}
\newtheorem{remark}[subsection]{Remark}
\newtheorem{remarks}[subsection]{Remarks}

\numberwithin{equation}{subsection}

% Macros
%
\def\lim{\mathop{\mathrm{lim}}\nolimits}
\def\colim{\mathop{\mathrm{colim}}\nolimits}
\def\Spec{\mathop{\mathrm{Spec}}}
\def\Hom{\mathop{\mathrm{Hom}}\nolimits}
\def\Ext{\mathop{\mathrm{Ext}}\nolimits}
\def\SheafHom{\mathop{\mathcal{H}\!\mathit{om}}\nolimits}
\def\SheafExt{\mathop{\mathcal{E}\!\mathit{xt}}\nolimits}
\def\Sch{\mathit{Sch}}
\def\Mor{\mathop{\mathrm{Mor}}\nolimits}
\def\Ob{\mathop{\mathrm{Ob}}\nolimits}
\def\Sh{\mathop{\mathit{Sh}}\nolimits}
\def\NL{\mathop{N\!L}\nolimits}
\def\CH{\mathop{\mathrm{CH}}\nolimits}
\def\proetale{{pro\text{-}\acute{e}tale}}
\def\etale{{\acute{e}tale}}
\def\QCoh{\mathit{QCoh}}
\def\Ker{\mathop{\mathrm{Ker}}}
\def\Im{\mathop{\mathrm{Im}}}
\def\Coker{\mathop{\mathrm{Coker}}}
\def\Coim{\mathop{\mathrm{Coim}}}

% Boxtimes
%
\DeclareMathSymbol{\boxtimes}{\mathbin}{AMSa}{"02}

%
% Macros for moduli stacks/spaces
%
\def\QCohstack{\mathcal{QC}\!\mathit{oh}}
\def\Cohstack{\mathcal{C}\!\mathit{oh}}
\def\Spacesstack{\mathcal{S}\!\mathit{paces}}
\def\Quotfunctor{\mathrm{Quot}}
\def\Hilbfunctor{\mathrm{Hilb}}
\def\Curvesstack{\mathcal{C}\!\mathit{urves}}
\def\Polarizedstack{\mathcal{P}\!\mathit{olarized}}
\def\Complexesstack{\mathcal{C}\!\mathit{omplexes}}
% \Pic is the operator that assigns to X its picard group, usage \Pic(X)
% \Picardstack_{X/B} denotes the Picard stack of X over B
% \Picardfunctor_{X/B} denotes the Picard functor of X over B
\def\Pic{\mathop{\mathrm{Pic}}\nolimits}
\def\Picardstack{\mathcal{P}\!\mathit{ic}}
\def\Picardfunctor{\mathrm{Pic}}
\def\Deformationcategory{\mathcal{D}\!\mathit{ef}}
\usepackage{cleveref}
\crefname{originallemma}{original lemma}{original lemmas}
\Crefname{originallemma}{Original lemma}{Original lemmas}
\makeatletter
\newlabel{homology-section-complexes}{{010V}{source}{}{sourcepointer-1}{}}
\newlabel{homology-section-complexes@cref}{{[originallemma][0][]010V}{[1][0][]source}{}{}{}}
\makeatother
\crefname{originallemma}{original lemma}{original lemmas}
\Crefname{originallemma}{Original lemma}{Original lemmas}
\makeatletter
\newlabel{homology-lemma-compose-homotopy}{{010W}{source}{}{sourcepointer-2}{}}
\newlabel{homology-lemma-compose-homotopy@cref}{{[originallemma][0][]010W}{[1][0][]source}{}{}{}}
\makeatother
\title{The cosimplicial Dold-Kan construction carries homotopic maps to homotopic maps}
\author{The Stacks Project Authors}
\date{stacks-project:simplicial.tex:local-snapshot-20261012; GNU FDL 1.2 or later}
\begin{document}
\maketitle
\noindent\textit{Editorial note.} The counted statement is the author's own lemma environment labelled lemma-homotopy-s-Q (source lines 5278-5288); the proof environment begins on the second line after it and is its own proof, closing at source line 5303.\par\medskip
\noindent\textit{Editorial note.} The statement and the proof invoke two author results of the same chapter by label; both are retained as context spans and the generator expands the context transitively (budget 800 lines for this card), while reference targets in other chapters that carry canonical project tags are auto-declared as non-load-bearing source pointers.\par\medskip
\noindent\textit{Editorial note.} The retained preamble is the chapter's own minimal preamble (sources/stacks-simplicial/preamble.tex, taken with its real line count), a verbatim excerpt of the native Xy-pic setup, the native theorem declarations with the mathematical macros, and the project's own no-op reference environments; no macro was invented or redefined.\par\medskip
\noindent\textit{Editorial scope.} One original result of the Stacks Project's Simplicial Methods chapter, counted once: the cosimplicial Dold-Kan construction carries homotopic maps to homotopic maps on s and on Q.\par\medskip
\noindent\textit{Allowed standard prerequisites (editorial).}\begin{itemize}
\item Cosimplicial objects and their cochain complexes. For a cosimplicial object U of an additive category, the associated cochain complex s(U) and, over an abelian category, the complex Q(U).. They are the objects of the statement (source lines 5283-5287).
\item Homotopy and the opposite category. Homotopic morphisms of cosimplicial objects correspond to homotopic morphisms in the opposite category via the contravariant functor A to A (the author's lemma-push-outs-simplicial-object-w-section, retained as context).. It transports the hypothesis to the opposite category (source lines 5291-5294).
\item The simplicial statement and the identification of the opposites. Homotopic maps of simplicial objects induce homotopic maps on s and on N (the author's lemma-homotopy-s-N, retained as context), together with the identities s(a') = (s(a))' and Q(U)' = N(U').. It supplies the homotopy in the opposite category and the two identifications (source lines 5295-5302).
\end{itemize}\medskip
\ifdefined\hypertarget\hypertarget{sourcepointer-1}{}\fi
\noindent\textit{Source reference.} Original lemma 010V: sources/stacks-chow/tags.txt, source lines 1341--1341. Cross-chapter pointer whose target carries the canonical tag 010V in the project tag list; declared so the reference resolves inside the assembled package, while every step the counted argument uses remains inside the retained spans.\par\medskip
\ifdefined\hypertarget\hypertarget{sourcepointer-2}{}\fi
\noindent\textit{Source reference.} Original lemma 010W: sources/stacks-chow/tags.txt, source lines 1342--1342. Cross-chapter pointer whose target carries the canonical tag 010W in the project tag list; declared so the reference resolves inside the assembled package, while every step the counted argument uses remains inside the retained spans.\par\medskip
\begin{lemma}
\label{lemma-push-outs-simplicial-object-w-section}
Let $f : X \to Y$ be a morphism of a category $\mathcal{C}$ with
pushouts. Assume there is a morphism $s : Y \to X$ with
$s \circ f = \text{id}_X$. Consider the cosimplicial object $U$ constructed in
Example \ref{example-push-outs-simplicial-object}
starting with $f$. The morphism $U \to U$ which in each degree
is the self map of $Y \amalg_X \ldots \amalg_X Y$
given by $f \circ s$ on each factor is homotopic to the identity on $U$.
In particular, $U$ is homotopy equivalent to the constant
cosimplicial object $X$.
\end{lemma}

\begin{proof}
This lemma is dual to
Lemma \ref{lemma-fibre-products-simplicial-object-w-section}.
Hence this lemma follows on applying
Lemma \ref{lemma-compare-homotopies}.
\end{proof}

\begin{lemma}
\label{lemma-homotopy-s-N}
Let $\mathcal{A}$ be an additive category. Let $a, b : U \to V$ be morphisms
of simplicial objects of $\mathcal{A}$. If $a$, $b$ are homotopic,
then $s(a), s(b) : s(U) \to s(V)$ are homotopic maps of chain complexes.
If $\mathcal{A}$ is abelian, then also $N(a), N(b) : N(U) \to N(V)$ are
homotopic maps of chain complexes.
\end{lemma}

\begin{proof}
We may choose a sequence $a = a_0, a_1, \ldots, a_n = b$ of morphisms from
$U$ to $V$ such that for each $i = 1, \ldots, n$ either there is a
homotopy from $a_i$ to $a_{i - 1}$ or there is a
homotopy from $a_{i - 1}$ to $a_i$. The calculation above
shows that in this case either
$s(a_i)$ is homotopic to $s(a_{i - 1})$ as a map of chain complexes
or $s(a_{i - 1})$ is homotopic to $s(a_i)$ as a map of chain complexes.
Of course, these things are equivalent and moreover being homotopic
is an equivalence relation on the set of maps of chain complexes, see
Homology, Section \ref{homology-section-complexes}.
This proves that $s(a)$ and $s(b)$ are homotopic as maps
of chain complexes.

\medskip\noindent
Next, we turn to $N(a)$ and $N(b)$. It follows from
Lemma \ref{lemma-decompose-associated-complexes} that
$N(a)$, $N(b)$ are compositions
$$
N(U) \to s(U) \to s(V) \to N(V)
$$
where we use $s(a)$, $s(b)$ in the middle. Hence the assertion
follows from
Homology, Lemma \ref{homology-lemma-compose-homotopy}.
\end{proof}

\begin{example}
\label{example-push-outs-simplicial-object}
Suppose that $X\to Y$ is a morphism of $C$ such that all
the pushouts $Y\amalg_X Y \amalg_X \ldots \amalg_X Y$ exist.
Then we set $U_n$ equal to the $(n + 1)$-fold pushout,
and we let $\varphi : [n] \to [m]$ correspond to the map
$$
(y \text{ in }i\text{th component})
\mapsto
(y \text{ in }\varphi(i)\text{th component})
$$
on ``coordinates''.
In other words, the map $U_1 = Y \amalg_X Y \to U_0 = Y$ is the
identity on each component.
The two maps $U_0 = Y \to U_1 = Y \amalg_X Y$ are the two coprojections.
\end{example}

\begin{lemma}
\label{lemma-compare-homotopies}
Let $\mathcal{C}$ be a category. Suppose that $U$ and $V$ are two
cosimplicial objects of $\mathcal{C}$. Let $a, b : U \to V$ be morphisms
of cosimplicial objects. Recall that $U$, $V$ correspond
to simplicial objects $U'$, $V'$ of $\mathcal{C}^{opp}$.
Moreover $a, b$ correspond to morphisms $a', b' : V' \to U'$.
The following are equivalent
\begin{enumerate}
\item There exists a homotopy $h = \{h_{n, \alpha}\}$ from
$a$ to $b$ as in Remark \ref{remark-homotopy-cosimplicial-better}.
\item There exists a homotopy $h = \{h_{n, i}\}$ from $a'$ to $b'$
as in Remark \ref{remark-homotopy-better}.
\end{enumerate}
Thus $a$ is homotopic to $b$ as in
Remark \ref{remark-homotopy-cosimplicial-better}
if and only if $a'$ is homotopic to $b'$ as in
Remark \ref{remark-homotopy-better}.
\end{lemma}

\begin{proof}
In case $\mathcal{C}$ has finite products, then $\mathcal{C}^{opp}$
has finite coproducts and we may use
Definitions \ref{definition-homotopy-cosimplicial}
and \ref{definition-homotopy} instead of
Remarks \ref{remark-homotopy-cosimplicial-better} and
\ref{remark-homotopy-better}.
In this case $h : U \to \Hom(\Delta[1], V)$
is the same as a morphism $h' : \Hom(\Delta[1], V)' \to U'$.
Since products and coproducts get switched too, it is immediate
that $(\Hom(\Delta[1], V))' = V' \times \Delta[1]$.
Moreover, the ``primed'' version of the morphisms
$e_0, e_1 : \Hom(\Delta[1], V) \to V$ are the morphisms
$e_0, e_1 : V' \to \Delta[1] \times V$.
Thus $e_0 \circ h = a$ translates into $h' \circ e_0 = a'$ and
similarly $e_1 \circ h = b$ translates into $h' \circ e_1 = b'$.
This proves the lemma in this case.

\medskip\noindent
In the general case, one needs to translate the relations given
by (\ref{equation-property-homotopy-cosimplicial}) into the relations
given in Lemma \ref{lemma-relations-homotopy}. We omit the details.

\medskip\noindent
The final assertion is formal from the equivalence of (1) and (2).
\end{proof}

\begin{lemma}
\label{lemma-fibre-products-simplicial-object-w-section}
Let $f : Y \to X$ be a morphism of a category $\mathcal{C}$ with
fibre products. Assume $f$ has a section $s$. Consider the
simplicial object $U$ constructed in
Example \ref{example-fibre-products-simplicial-object}
starting with $f$. The morphism $U \to U$ which in each degree
is the self map $(s \circ f)^{n + 1}$ of $Y \times_X \ldots \times_X Y$
given by $s \circ f$ on each factor is homotopic to the identity on $U$.
In particular, $U$ is homotopy equivalent to the constant
simplicial object $X$.
\end{lemma}

\begin{proof}
Set $g^0 = \text{id}_Y$ and $g^1 = s \circ f$.
We use the morphisms
\begin{eqnarray*}
Y \times_X \ldots \times_X Y \times \Mor([n], [1])
& \to &
Y \times_X \ldots \times_X Y \\
(y_0, \ldots, y_n) \times \alpha
& \mapsto &
(g^{\alpha(0)}(y_0), \ldots, g^{\alpha(n)}(y_n))
\end{eqnarray*}
where we use the functor of points point of view to define the maps.
Another way to say this is to say that
$h_{n, 0} = \text{id}$, $h_{n, n + 1} = (s \circ f)^{n + 1}$ and
$h_{n, i} = \text{id}_Y^{i + 1} \times (s \circ f)^{n + 1 - i}$.
We leave it to the reader to show that these satisfy the relations
of Lemma \ref{lemma-relations-homotopy}. Hence they define
the desired homotopy. See also Remark \ref{remark-homotopy-better}
which shows that
we do not need to assume anything else on the category $\mathcal{C}$.
\end{proof}

\begin{lemma}
\label{lemma-decompose-associated-complexes}
Let $\mathcal{A}$ be an abelian category.
Let $U$ be a simplicial object of $\mathcal{A}$.
The canonical morphism of chain complexes
$N(U) \to s(U)$ is split. In fact,
$$
s(U) = N(U) \oplus D(U)
$$
for some complex $D(U)$. The construction $U \mapsto D(U)$
is functorial.
\end{lemma}

\begin{proof}
Define $D(U)_n$ to be the image of
$$
\bigoplus\nolimits_{\varphi : [n] \to [m]\text{ surjective}, \ m < n} N(U_m)
\xrightarrow{\bigoplus U(\varphi)}
U_n
$$
which is a subobject of $U_n$ complementary to
$N(U_n)$ according to Lemma \ref{lemma-splitting-abelian-category} and
Definition \ref{definition-split}. We show that
$D(U)$ is a subcomplex. Pick a surjective
map $\varphi : [n] \to [m]$ with $m < n$ and consider
the composition
$$
N(U_m) \xrightarrow{U(\varphi)} U_n \xrightarrow{d_n} U_{n - 1}.
$$
This composition is the sum of the maps
$$
N(U_m) \xrightarrow{U(\varphi \circ \delta^n_i)} U_{n - 1}
$$
with sign $(-1)^i$, $i = 0, \ldots, n$.

\medskip\noindent
First we will prove by ascending induction on $m$,
$0 \leq m < n - 1$ that all the maps $U(\varphi \circ \delta^n_i)$
map $N(U_m)$ into $D(U)_{n - 1}$. (The case $m = n - 1$ is treated below.)
Whenever the map $\varphi \circ \delta^n_i : [n - 1] \to [m]$
is surjective then the image of $N(U_m)$ under $U(\varphi \circ \delta^n_i)$
is contained in $D(U)_{n - 1}$ by definition.
If $\varphi \circ \delta^n_i : [n - 1] \to [m]$ is not surjective,
set $j = \varphi(i)$ and observe that $i$ is the unique
index whose image under $\varphi$ is $j$. We may write
$\varphi \circ \delta^n_i = \delta^m_j \circ \psi \circ \delta^n_i$
for some $\psi : [n - 1] \to [m - 1]$. Hence
$U(\varphi \circ \delta^n_i) = U(\psi \circ \delta^n_i) \circ d^m_j$ which
is zero on $N(U_m)$ unless $j = m$. If $j = m$, then
$d^m_m(N(U_m)) \subset N(U_{m - 1})$ and hence
$U(\varphi \circ \delta^n_i)(N(U_m)) \subset
U(\psi \circ \delta^n_i)(N(U_{m - 1}))$ and we win
by induction hypothesis.

\medskip\noindent
To finish proving that $D(U)$ is a subcomplex
we still have to deal with the composition
$$
N(U_m) \xrightarrow{U(\varphi)} U_n \xrightarrow{d_n} U_{n - 1}.
$$
in case $m =  n - 1$. In this case $\varphi = \sigma^{n - 1}_j$
for some $0 \leq j \leq n - 1$ and $U(\varphi) = s^{n - 1}_j$.
Thus the composition is given by the sum
$$
\sum (-1)^i d^n_i \circ s^{n - 1}_j
$$
Recall from Remark \ref{remark-relations} that
$d^n_j \circ s^{n - 1}_j = d^n_{j + 1} \circ s^{n - 1}_j = \text{id}$
and these drop out because the corresponding terms have opposite signs.
The map $d^n_n \circ s^{n - 1}_j$, if $j < n - 1$, is equal to
$s^{n - 2}_j \circ d^{n - 1}_{n - 1}$. Since
$d^{n - 1}_{n - 1}$ maps $N(U_{n - 1})$ into $N(U_{n - 2})$,
we see that the image $d^n_n ( s^{n - 1}_j (N(U_{n - 1}))$
is contained in $s^{n - 2}_j(N(U_{n - 2}))$ which
is contained in $D(U_{n - 1})$ by definition. For all
other combinations of $(i, j)$ we have
either $d^n_i \circ s^{n - 1}_j = s^{n - 2}_{j - 1} \circ d^{n - 1}_i$
(if $i < j$), or
$d^n_i \circ s^{n - 1}_j = s^{n - 2}_j \circ d^{n - 1}_{i - 1}$
(if $n > i > j + 1$) and in these cases the map is zero because
of the definition of $N(U_{n - 1})$.
\end{proof}

\begin{definition}
\label{definition-homotopy}
Let $\mathcal{C}$ be a category having finite coproducts.
Suppose that $U$ and $V$ are two simplicial objects of $\mathcal{C}$.
Let $a, b : U \to V$ be two morphisms.
\begin{enumerate}
\item We say a morphism
$$
h : U \times \Delta[1] \longrightarrow V
$$
is a {\it homotopy from $a$ to $b$} if $a = h \circ e_0$ and
$b = h \circ e_1$.
\item We say the morphisms $a$ and $b$ are {\it homotopic} or are
{\it in the same homotopy class}
if there exists a sequence of morphisms $a = a_0, a_1, \ldots, a_n = b$
from $U$ to $V$ such that for each $i = 1, \ldots, n$ there either exists
a homotopy from $a_{i - 1}$ to $a_i$ or there exists a homotopy
from $a_i$ to $a_{i - 1}$.
\end{enumerate}
\end{definition}

\begin{definition}
\label{definition-homotopy-cosimplicial}
Let $\mathcal{C}$ be a category having finite products.
Let $U$ and $V$ be two cosimplicial objects of $\mathcal{C}$.
Let $a, b : U \to V$ be two morphisms of cosimplicial objects
of $\mathcal{C}$.
\begin{enumerate}
\item We say a morphism
$$
h : U \longrightarrow \Hom(\Delta[1], V)
$$
such that $a = e_0 \circ h$ and $b = e_1 \circ h$ is a
{\it homotopy from $a$ to $b$}.
\item We say $a$ and $b$ are {\it homotopic} or are
{\it in the same homotopy class} if there exists a sequence
$a = a_0, a_1, \ldots, a_n = b$ of morphisms from $U$ to $V$
such that for each $i = 1, \ldots, n$ there either exists a
homotopy from $a_i$ to $a_{i - 1}$ or there exists a homotopy
from $a_{i - 1}$ to $a_i$.
\end{enumerate}
\end{definition}

\begin{equation}
\label{equation-property-homotopy-cosimplicial}
h_{m, \alpha} \circ U(f) = V(f) \circ h_{n, \alpha \circ f}
\end{equation}

\begin{lemma}
\label{lemma-relations-homotopy}
In the situation above, we have the following relations:
\begin{enumerate}
\item We have $h_{n, 0} = b_n$ and $h_{n, n + 1} = a_n$.
\item We have $d^n_j \circ h_{n, i} = h_{n - 1, i - 1} \circ d^n_j$
for $i > j$.
\item We have $d^n_j \circ h_{n, i} = h_{n - 1, i} \circ d^n_j$
for $i \leq j$.
\item We have $s^n_j \circ h_{n, i} = h_{n + 1, i + 1} \circ s^n_j$
for $i > j$.
\item We have $s^n_j \circ h_{n, i} = h_{n + 1, i} \circ s^n_j$
for $i \leq j$.
\end{enumerate}
Conversely, given a system of maps $h_{n, i}$ satisfying the
properties listed above, then these define a morphism
$h$ which is a homotopy from $a$ to $b$.
\end{lemma}

\begin{proof}
Omitted. You can prove the last statement using the fact,
see Lemma \ref{lemma-face-degeneracy-category} that
to give a morphism of simplicial objects is the
same as giving a sequence of morphisms $h_n$ commuting
with all $d^n_j$ and $s^n_j$.
\end{proof}

\begin{remark}
\label{remark-homotopy-better}
Let $\mathcal{C}$ be any category (no assumptions whatsoever). Let
$U$ and $V$ be simplicial objects of $\mathcal{C}$. Let $a, b : U \to V$
be morphisms of simplicial objects of $\mathcal{C}$. A
{\it homotopy from $a$ to $b$} is given by
morphisms\footnote{In the literature, often the maps
$h_{n + 1, i} \circ s_i : U_n \to V_{n + 1}$ are used instead
of the maps $h_{n, i}$. Of course the relations these maps satisfy
are different from the ones in Lemma \ref{lemma-relations-homotopy}.}
$h_{n, i} : U_n \to V_n$, for $n \geq 0$, $i = 0, \ldots, n + 1$
satisfying the relations of Lemma \ref{lemma-relations-homotopy}.
As in Definition \ref{definition-homotopy} we say the morphisms $a$ and $b$
are {\it homotopic} if there exists a sequence of morphisms
$a = a_0, a_1, \ldots, a_n = b$ from $U$ to $V$ such that for each
$i = 1, \ldots, n$ there either exists a homotopy from $a_{i - 1}$ to $a_i$
or there exists a homotopy from $a_i$ to $a_{i - 1}$.
Clearly, if $F : \mathcal{C} \to \mathcal{C}'$ is any functor
and $\{h_{n, i}\}$ is a homotopy from $a$ to $b$, then
$\{F(h_{n, i})\}$ is a homotopy from $F(a)$ to $F(b)$.
Similarly, if $a$ and $b$ are homotopic, then $F(a)$ and $F(b)$
are homotopic.
Since the lemma says that the newer notion is the same
as the old one in case finite coproduct exist, we deduce
in particular that functors preserve the original notion
whenever both categories have finite coproducts.
\end{remark}

\begin{remark}
\label{remark-homotopy-cosimplicial-better}
Let $\mathcal{C}$ be any category (no assumptions whatsoever). Let
$U$ and $V$ be cosimplicial objects of $\mathcal{C}$. Let $a, b : U \to V$
be morphisms of cosimplicial objects of $\mathcal{C}$. A
{\it homotopy from $a$ to $b$} is given by morphisms
$h_{n, \alpha} : U_n \to V_n$, for $n \geq 0$, $\alpha \in \Delta[1]_n$
satisfying (\ref{equation-property-homotopy-cosimplicial})
for all morphisms $f$ of $\Delta$ and such that
$a_n = h_{n, 0 : [n] \to [1]}$ and $b_n = h_{n, 1 : [n] \to [1]}$
for all $n \geq 0$.
As in Definition \ref{definition-homotopy-cosimplicial}
we say the morphisms $a$ and $b$
are {\it homotopic} if there exists a sequence of morphisms
$a = a_0, a_1, \ldots, a_n = b$ from $U$ to $V$ such that for each
$i = 1, \ldots, n$ there either exists a homotopy from $a_{i - 1}$ to $a_i$
or there exists a homotopy from $a_i$ to $a_{i - 1}$.
Clearly, if $F : \mathcal{C} \to \mathcal{C}'$ is any functor
and $\{h_{n, i}\}$ is a homotopy from $a$ to $b$, then
$\{F(h_{n, i})\}$ is a homotopy from $F(a)$ to $F(b)$.
Similarly, if $a$ and $b$ are homotopic, then $F(a)$ and $F(b)$
are homotopic. This new notion is the same
as the old one in case finite products exist. We deduce
in particular that functors preserve the original notion
whenever both categories have finite products.
\end{remark}

\begin{example}
\label{example-fibre-products-simplicial-object}
Suppose that $Y\to X$ is a morphism of $\mathcal{C}$ such that all
the fibred products $Y \times_X Y \times_X \ldots \times_X Y$ exist.
Then we set $U_n$ equal to the $(n + 1)$-fold fibre product,
and we let $\varphi : [n] \to [m]$ correspond to the map
(on ``coordinates'')
$(y_0, \ldots, y_m) \mapsto (y_{\varphi(0)}, \ldots, y_{\varphi(n)})$.
In other words, the map $U_0 = Y \to U_1 = Y \times_X Y$ is the
diagonal map. The two maps $U_1 = Y \times_X Y \to U_0 = Y$ are the
projection maps.
\end{example}

\begin{definition}
\label{definition-split}
Let $\mathcal{C}$ be a category which admits finite nonempty coproducts.
We say a simplicial object $U$ of $\mathcal{C}$ is {\it split}
if there exist subobjects $N(U_m)$ of $U_m$, $m \geq 0$
with the property that
\begin{equation}
\label{equation-splitting}
\coprod\nolimits_{\varphi : [n] \to [m]\text{ surjective}}
N(U_m)
\longrightarrow
U_n
\end{equation}
is an isomorphism for all $n \geq 0$. If $U$ is an $r$-truncated
simplicial object of $\mathcal{C}$ then we say $U$ is {\it split}
if there exist subobjects $N(U_m)$ of $U_m$, $r \geq m \geq 0$
with the property that (\ref{equation-splitting})
is an isomorphism for $r \geq n \geq 0$.
\end{definition}

\begin{lemma}
\label{lemma-splitting-abelian-category}
Let $\mathcal{A}$ be an abelian category.
Let $U$ be a simplicial object in $\mathcal{A}$.
Then $U$ has a splitting obtained by taking $N(U_0) = U_0$ and
for $m \geq 1$ taking
$$
N(U_m) = \bigcap\nolimits_{i = 0}^{m - 1} \Ker(d^m_i).
$$
Moreover, this splitting is functorial on the category of
simplicial objects of $\mathcal{A}$.
\end{lemma}

\begin{proof}
For any object $A$ of $\mathcal{A}$ we obtain
a simplicial abelian group $\Mor_\mathcal{A}(A, U)$.
Each of these are canonically split by Lemma
\ref{lemma-splitting-simplicial-groups}. Moreover,
$$
N(\Mor_\mathcal{A}(A, U_m)) =
\bigcap\nolimits_{i = 0}^{m - 1} \Ker(d^m_i) =
\Mor_\mathcal{A}(A, N(U_m)).
$$
Hence we see that the morphism (\ref{equation-splitting})
becomes an isomorphism after applying the functor
$\Mor_\mathcal{A}(A, -)$ for any object of $\mathcal{A}$.
Hence it is an isomorphism by the Yoneda lemma.
\end{proof}

\begin{remark}
\label{remark-relations}
By abuse of notation we sometimes write $d_i : U_n \to U_{n - 1}$
instead of $d^n_i$, and similarly for $s_i : U_n \to U_{n + 1}$.
The relations among the morphisms $d^n_i$ and $s^n_i$
may be expressed as follows:
\begin{enumerate}
\item If $i < j$, then $d_i \circ d_j = d_{j - 1} \circ d_i$.
\item If $i < j$, then $d_i \circ s_j = s_{j - 1} \circ d_i$.
\item We have $\text{id} = d_j \circ s_j = d_{j + 1} \circ s_j$.
\item If $i > j + 1$, then $d_i \circ s_j = s_j \circ d_{i - 1}$.
\item If $i \leq j$, then $s_i \circ s_j = s_{j + 1} \circ s_i$.
\end{enumerate}
This means that whenever the compositions on both the left and the
right are defined then the corresponding equality should hold.
\end{remark}

\begin{lemma}
\label{lemma-face-degeneracy-category}
The category $\Delta$ is the universal category
with objects $[n]$, $n \geq 0$ and morphisms
$\delta^n_j$ and $\sigma^n_j$ such that (a) every morphism is
a composition of these morphisms, (b) the relations
listed in Lemma \ref{lemma-relations-face-degeneracy} are satisfied,
and (c) any relation among the morphisms is a consequence of
those relations.
\end{lemma}

\begin{proof}
Omitted.
\end{proof}

\begin{lemma}
\label{lemma-splitting-simplicial-groups}
Let $U$ be a simplicial abelian group.
Then $U$ has a splitting obtained by taking $N(U_0) = U_0$ and
for $m \geq 1$ taking
$$
N(U_m) = \bigcap\nolimits_{i = 0}^{m - 1} \Ker(d^m_i).
$$
Moreover, this splitting is functorial on the category
of simplicial abelian groups.
\end{lemma}

\begin{proof}
By induction on $n$ we will show that the choice of $N(U_m)$
in the lemma guarantees that (\ref{equation-splitting}) is
an isomorphism for $m \leq n$. This is clear for $n = 0$.
In the rest of this proof we are going to
drop the superscripts from the maps $d_i$ and $s_i$ in order
to improve readability. We will also repeatedly use the relations
from Remark \ref{remark-relations}.

\medskip\noindent
First we make a general remark.
For $0 \leq i \leq m$ and $z \in U_m$ we have
$d_i(s_i(z)) = z$. Hence we can write
any $x \in U_{m + 1}$ uniquely as
$x = x' + x''$ with $d_i(x') = 0$
and $x'' \in \Im(s_i)$
by taking $x' = (x - s_i(d_i(x)))$ and
$x'' = s_i(d_i(x))$. Moreover, the element
$z \in U_m$ such that $x'' = s_i(z)$
is unique because $s_i$ is injective.

\medskip\noindent
Here is a procedure for decomposing any $x \in U_{n + 1}$.
First, write $x = x_0 + s_0(z_0)$ with $d_0(x_0) = 0$.
Next, write $x_0 = x_1 + s_1(z_1)$ with
$d_1(x_1) = 0$. Continue like this to get
\begin{eqnarray*}
x & = & x_0 + s_0(z_0), \\
x_0 & = & x_1 + s_1(z_1), \\
x_1 & = & x_2 + s_2(z_2), \\
\ldots & \ldots & \ldots \\
x_{n - 1} & = & x_n + s_n(z_n)
\end{eqnarray*}
where $d_i(x_i) = 0$ for all $i = n, \ldots, 0$.
By our general remark above all of the $x_i$
and $z_i$ are determined uniquely by $x$.
We claim that
$x_i \in
\Ker(d_0) \cap
\Ker(d_1) \cap
\ldots \cap
\Ker(d_i)$
and
$z_i \in
\Ker(d_0) \cap
\ldots \cap
\Ker(d_{i - 1})$
for $i = n, \ldots, 0$.
Here and in the following
an empty intersection of kernels indicates
the whole space; i.e.,
the notation
$z_0 \in \Ker(d_0) \cap
\ldots \cap
\Ker(d_{i - 1})$
when $i = 0$ means $z_0 \in U_n$ with no restriction.

\medskip\noindent
We prove this by ascending induction on $i$.
It is clear for $i = 0$ by construction of $x_0$ and $z_0$.
Let us prove it for $0 < i \leq n$ assuming the result for $i - 1$.
First of all we have $d_i(x_i) = 0$ by construction.
So pick a $j$ with $0 \leq j < i$. We have
$d_j(x_{i - 1}) = 0$ by induction. Hence
$$
0 = d_j(x_{i - 1})
= d_j(x_i) + d_j(s_i(z_i))
= d_j(x_i) + s_{i - 1}(d_j(z_i)).
$$
The last equality by the relations of Remark \ref{remark-relations}.
These relations also imply that
$d_{i - 1}(d_j(x_i)) = d_j(d_i(x_i)) = 0$
because $d_i(x_i)= 0$ by construction.
Then the uniqueness in the general remark above shows the equality
$0 = x' + x'' = d_j(x_i) + s_{i - 1}(d_j(z_i))$
can only hold if both terms are zero. We conclude that
$d_j(x_i) = 0$ and by injectivity of $s_{i - 1}$ we also
conclude that $d_j(z_i) = 0$. This proves the claim.

\medskip\noindent
The claim implies we can uniquely write
$$
x = s_0(z_0) + s_1(z_1) + \ldots + s_n(z_n) + x_n
$$
with $x_n \in N(U_{n + 1})$ and
$z_i \in \Ker(d_0) \cap \ldots \cap \Ker(d_{i - 1})$.
We can reformulate this as saying that we have found a direct
sum decomposition
$$
U_{n + 1}
=
N(U_{n + 1})
\oplus
\bigoplus\nolimits_{i = 0}^{i = n}
s_i\Big(\Ker(d_0) \cap \ldots \cap \Ker(d_{i - 1})\Big)
$$
with the property that
$$
\Ker(d_0) \cap \ldots \cap \Ker(d_j)
=
N(U_{n + 1}) \oplus
\bigoplus\nolimits_{i = j + 1}^{i = n}
s_i\Big(\Ker(d_0) \cap \ldots \cap \Ker(d_{i - 1})\Big)
$$
for $j = 0, \ldots, n$.
The result follows from this statement as follows.
Each of the $z_i$ in the expression for $x$
can be written uniquely as
$$
z_i = s_i(z'_{i, i}) + \ldots + s_{n - 1}(z'_{i, n - 1}) + z_{i, 0}
$$
with $z_{i, 0} \in N(U_n)$ and
$z'_{i, j} \in \Ker(d_0) \cap \ldots \cap \Ker(d_{j - 1})$.
The first few steps in the decomposition of $z_i$ are zero because
$z_i$ already is in the kernel of $d_0, \ldots, d_i$.
This in turn uniquely gives
$$
x = x_n + s_0(z_{0, 0}) + s_1(z_{1, 0}) + \ldots + s_n(z_{n, 0}) +
\sum\nolimits_{0 \leq i \leq j \leq n - 1} s_i(s_j(z'_{i, j})).
$$
Continuing in this fashion we see that we in the end obtain
a decomposition of $x$ as a sum of terms
of the form
$$
s_{i_1} s_{i_2} \ldots s_{i_k} (z)
$$
with $0 \leq i_1 \leq i_2 \leq \ldots \leq i_k \leq n - k + 1$ and
$z \in N(U_{n + 1 - k})$. This is exactly the required
decomposition, because any surjective map $[n + 1] \to [n + 1 - k]$
can be uniquely expressed in the form
$$
\sigma^{n - k}_{i_k} \ldots \sigma^{n - 1}_{i_2} \sigma^n_{i_1}
$$
with $0 \leq i_1 \leq i_2 \leq \ldots \leq i_k \leq n - k + 1$.
\end{proof}

\begin{lemma}
\label{lemma-relations-face-degeneracy}
The morphisms $\delta^n_j$ and $\sigma^n_j$ satisfy the following relations.
\begin{enumerate}
\item If $0 \leq i < j \leq n + 1$, then
$\delta^{n + 1}_j \circ \delta^n_i =
\delta^{n + 1}_i \circ \delta^n_{j - 1}$.
In other words the diagram
$$
\xymatrix{
& [n] \ar[rd]^{\delta^{n + 1}_j} & \\
[n - 1] \ar[ru]^{\delta^n_i} \ar[rd]_{\delta^n_{j - 1}} & &
[n + 1] \\
& [n] \ar[ru]_{\delta^{n + 1}_i} &
}
$$
commutes.
\item If $0 \leq i < j \leq n - 1$, then
$\sigma^{n - 1}_j \circ \delta^n_i =
\delta^{n - 1}_i \circ \sigma^{n - 2}_{j - 1}$.
In other words the diagram
$$
\xymatrix{
& [n] \ar[rd]^{\sigma^{n - 1}_j} & \\
[n - 1] \ar[ru]^{\delta^n_i} \ar[rd]_{\sigma^{n - 2}_{j - 1}} & &
[n - 1] \\
& [n - 2] \ar[ru]_{\delta^{n - 1}_i} &
}
$$
commutes.
\item If $0 \leq j \leq n - 1$, then
$\sigma^{n - 1}_j \circ \delta^n_j = \text{id}_{[n - 1]}$
and
$\sigma^{n - 1}_j \circ \delta^n_{j + 1} = \text{id}_{[n - 1]}$.
In other words the diagram
$$
\xymatrix{
& [n] \ar[rd]^{\sigma^{n - 1}_j} & \\
[n - 1]
\ar[ru]^{\delta^n_j}
\ar[rd]_{\delta^n_{j + 1}}
\ar[rr]^{\text{id}_{[n - 1]}} & & [n - 1] \\
& [n] \ar[ru]_{\sigma^{n - 1}_j} &
}
$$
commutes.
\item If $0 < j + 1 < i \leq n$, then
$\sigma^{n - 1}_j \circ \delta^n_i =
\delta^{n - 1}_{i - 1} \circ \sigma^{n - 2}_j$.
In other words the diagram
$$
\xymatrix{
& [n] \ar[rd]^{\sigma^{n - 1}_j} & \\
[n - 1] \ar[ru]^{\delta^n_i} \ar[rd]_{\sigma^{n - 2}_j} & &
[n - 1] \\
& [n - 2] \ar[ru]_{\delta^{n - 1}_{i - 1}} &
}
$$
commutes.
\item If $0 \leq i \leq j \leq n - 1$, then
$\sigma^{n - 1}_j \circ \sigma^n_i =
\sigma^{n - 1}_i \circ \sigma^n_{j + 1}$.
In other words the diagram
$$
\xymatrix{
& [n] \ar[rd]^{\sigma^{n - 1}_j} & \\
[n + 1] \ar[ru]^{\sigma^n_i} \ar[rd]_{\sigma^n_{j + 1}} & &
[n - 1] \\
& [n] \ar[ru]_{\sigma^{n - 1}_i} &
}
$$
commutes.
\end{enumerate}
\end{lemma}

\begin{proof}
Omitted.
\end{proof}

\begin{lemma}
\label{lemma-homotopy-s-Q}
\begin{slogan}
The (cosimplicial) Dold-Kan functor carries homotopic maps to homotopic maps.
\end{slogan}
Let $\mathcal{A}$ be an additive category. Let $a, b : U \to V$ be morphisms
of cosimplicial objects of $\mathcal{A}$. If $a$, $b$ are homotopic,
then $s(a), s(b) : s(U) \to s(V)$ are homotopic maps of cochain complexes.
If in addition $\mathcal{A}$ is abelian, then $Q(a), Q(b) : Q(U) \to Q(V)$
are homotopic maps of cochain complexes.
\end{lemma}

\begin{proof}
Let $(-)' : \mathcal{A} \to \mathcal{A}^{opp}$
be the contravariant functor $A \mapsto A$. By
Lemma \ref{lemma-push-outs-simplicial-object-w-section}
the maps $a'$ and $b'$ are homotopic.
By Lemma \ref{lemma-homotopy-s-N} we see that
$s(a')$ and $s(b')$ are homotopic maps of chain
complexes. Since $s(a') = (s(a))'$ and
$s(b') = (s(b))'$ we conclude that also
$s(a)$ and $s(b)$ are homotopic by applying the additive
contravariant functor $(-)'' : \mathcal{A}^{opp} \to \mathcal{A}$.
The result for the $Q$-complexes follows in the same manner
using that $Q(U)' = N(U')$.
\end{proof}

\end{document}
